\section{Introduction}

Some exoplanet climate studies hold a planet's air close to Earth's and vary its rotation, gravity or star.
\textcite{yang2014strong} and \textcite{way2018climates} slowed the rotation of planets with nitrogen-dominated
air and followed their cloud and temperature as the day lengthened, \textcite{thomson2019effects} varied the
gravity of a terrestrial planet in moist and dry general circulation models, and \textcite{boutle2017exploring}
placed a nitrogen--oxygen atmosphere on Proxima Centauri~b, under a cooler star. For such work, by the account of
THOR's developers, most general circulation model (GCM) studies of exoplanet atmospheres adapt codes developed for
Earth or for solar-system planets \autocite{deitrick2020thor}: ROCKE-3D was developed from the GISS ModelE2, ExoCAM
from NCAR's Community Atmosphere Model, and ExoPlaSim from the Planet Simulator, PlaSim
\autocite{way2017resolving,wolf2022exocam,paradise2022exoplasim}. Convection and wave models built for Earth have
been run off Earth as well: SAM in cloud-permitting runs on tidally locked planets, CM1 in pure-steam air, and the
coastal wave model SWAN on the ancient seas of Mars \autocite{yang2023cloud,tan2021convection,banfield2015winds}.
Other models built for Earth have been joined to GCMs. On Titan, \textcite{rafkin2022interaction} started a
cloud-resolving model from a GCM's unstable profiles, \textcite{battalio2022interaction} tuned the GCM's convection
scheme with the cloud-resolving runs and observations, and a GCM's winds have driven a wave model over the polar
lakes \autocite{detelich2026modeling}. The LMD's GCM physics, for its part, runs on WRF's core, in a mesoscale
model on Mars and a convection-resolving model over Proxima~b, so that the global and resolved calculations share
their physics \autocite{spiga2009new,lefevre2021convection}.

Constructed worlds can depart further from Earth, for the bodies of the Solar System that keep atmospheres are those
whose escape velocity is large against the light they receive: the Moon held the volcanic air of its youth, at
most about 1~kPa, for an estimated 70 million years or more \autocite{needham2017lunar}, while Titan, far from the Sun and at
94~K, keeps a dense atmosphere on a body smaller than Mars. Beyond the Solar System, the known exoplanets cluster
near the same divide, extended to them \autocite{zahnle2017cosmic}. An engineered world, by contrast, can hold
nitrogen--oxygen air where a natural one would lose it. \textcite{mckay1991making} asked how Mars could be made
habitable, and the Open Moon, a breathable Moon designed in the Terluna project \autocite{goldman2026terluna}, is
a design of that kind for the Moon: it is to keep 1.2~atm of nitrogen and oxygen under a sixth of Earth's gravity
behind a shield in space that removes the ultraviolet heating its upper air. By the calculations of atmospheric
escape made for the design, the Moon's exobase stays below 200~K and the air loses 0.6--1.4~kg\,s$^{-1}$ over a solar cycle, losses that would take 60
to 150 billion years to remove it.\claim{S15} Such worlds combine nitrogen--oxygen air with a gravity, a size and
a day that no known natural world combines with it.

Modelling such worlds means adapting codes like those named above, in which the developers who ported Earth
models to other planets found values fixed for Earth. PlanetWRF's developers, finding WRF written to apply to
Earth alone, redefined its second as 1/86,400 of the planet's rotation period and removed hard-wired Earth
settings from its physics, and built TitanWRF on the result
\autocite{richardson2007planetwrf,newman2011stratospheric}. ROCKE-3D's developers found constants set for the lengths of Earth's day and year in its first release and, in its second, removed or generalised a set of
Earth-centric options, measuring their effects in a sequence of runs
\autocite{way2017resolving,tsigaridis2025rocke3d}; ExoPlaSim's description records that PlaSim's solar day had
to be taken out of its unit conversions before slow rotators could be run, and that its near-infrared water
absorptivity is tuned to a solar spectrum \autocite{paradise2022exoplasim}. Each team removed Earth values
that mattered for the worlds it set out to model, but others could stay: ROCKE-3D's developers report that the
changes from the Earth configuration in their first release went without systematic study \autocite{tsigaridis2025rocke3d},
and THOR's developers, who wrote their model for other planets from the start, describe planet-specific tunings
as often built into adapted codes and hard to find \autocite{deitrick2020thor}. On a Gaean world, then, a run can
compute the familiar phases of water in nitrogen--oxygen air while its code still contains a constant set for Earth's gravity,
or a formula that counts rain per Earth day, and because its output is plausible, the retained value can be hard
to detect. Plausible results can still disagree: four GCMs given the same moist cases of TRAPPIST-1e all simulate a
temperate climate, yet differ in global-mean surface temperature by 14~K with nitrogen-dominated air and 24~K with carbon
dioxide, a spread their authors relate to the models' parameterization choices \autocite{sergeev2022thai}. In an
ensemble linking global climate to resolved storms and coastal waves, moreover, models built separately and tuned
to Earth in different ways pass forcing to one another, so the retained Earth values affect both the individual
calculations and the fields exchanged between them. Developing such an ensemble for Gaean exoplanets
and constructed worlds therefore requires identifying which changes follow from general laws and which must be calculated or calibrated for each world.

Toward such an ensemble, the adaptations and couplings described here began with the Open Moon as their test
world. Its air is close to Earth's in composition and surface pressure, while four ratios on which the models'
Earth values depend are far from Earth's at once: the column over each square metre holds 7.25 times Earth's air,\claim{S02} the
scale height of about 53~km is 23 times Earth's in proportion to the radius,\claim{S03} gravity is a sixth of
Earth's, and the solar day lasts 29.53 Earth days. In a run of a natural planet that departs from Earth along one of these ratios, a slow rotator or a smaller world, only the Earth values that depend on that ratio err, whereas
on the test world the values of all four err together, by amounts large enough to measure. PlaSim, for example,
fixes the horizon air mass at Earth's 35, where the Moon's well-mixed air has about 7. The design's seas and lakes
also set the test world's orientation: their uneven load turns the nearly spherical Moon about 8\textdegree\ beneath
its spin axis while its crust responds elastically, a case of the true polar wander that mass loads drive
\autocite{matsuyama2014planetary}, and the geography and tilt under which the models run follow from that
turn.\claim{S17}

The global climate is computed by ExoPlaSim~3.4.2 \autocite{paradise2022exoplasim,paradise2025exoplasim}, PlaSim
\autocite{fraedrich2005planet} extended for other worlds (the appendix gives each code's configuration), which runs
decades of the test world's climate in a few hours on four processor cores, the speed for which it was
chosen.\claim{S14} On fifteen levels laid out for the test world's scale height, the GCM parameterizes its convection
with a simplified Betts--Miller scheme \autocite{frierson2007dynamics}, computes its radiation with broadband fits to
Earth's columns and diagnoses its clouds from formulas fitted to Earth.\claim{R20}\claim{R21} Slices of the GCM's
climate force CM1 release~22 \autocite{bryan2002benchmark,bryan2026cm1} on two-dimensional rings of columns around
latitude circles or tilted great circles and in doubly periodic boxes, where it resolves the convection the GCM
parameterizes and forms its own clouds and rain. In the boxes run to compare the two models, the GCM's heating
and moistening enter through the whole column, and CM1's temperature and moisture are held to the GCM's only above
the GCM's convection top, a height the forcing model computes on any world.\claim{R14} With the electrification of WRF-ELEC grafted in
\autocite{mansell2010simulated,fierro2013implementation,mansell2026wrfelec}, it also separates charge and makes
lightning. These rings and boxes were sized to fit the hardware at hand, one laptop with eight processor cores,
and chosen on a best guess at which features would be of interest.\claim{S16} The GCM's surface stress drives
SWAN~41.51 \autocite{booij1999third,swan4151user} over the seas, and
SWASH~12.01 \autocite{zijlema2011swash,swash1201user}, forced by SWAN's spectra, computes the waves at the shore.
Alongside these couplings, the radiation of both atmospheric models was tested against an in-house line-by-line
radiative--convective model, the reference model below, built on HITRAN lines as served in September 2026
\autocite{gordon2022hitran2020} and the MT\_CKD continuum \autocite{mlawer2012development}: the GCM's radiation
was calibrated to it, and CM1's checked against it.

Adapting these codes to the test world meant finding two things in them: Earth values, each a value or form in
a model that agrees with Earth because it was set or fitted there, and defects, code that contradicts its own
inputs, units or intent. In the four codes the Earth values have four forms: dimensional constants written into
the code, units of time, geometry, and empirical fits and simplifications used beyond the range they were made
for. On the paths the lunar runs take, the released codes also contain five defects. Three of them produce
little or no error at Earth's values yet, by ExoPlaSim's own arithmetic, bias its runs that meet their
conditions, runs of slowly rotating and tidally locked planets among them, while the other two err on Earth as well:
CM1's simple boundary layer multiplies its mixing by the air's density in kilograms per cubic metre, and SWAN reads
a parameter it has left unset.\claim{H34} Several of these values and defects leave a run plausible on its own; they
were found by setting one model against another or against its own surface fluxes, or by reading the
code.\claim{H34}\claim{H38} Of the repairs made to these values
and defects, most belong to three classes, which differ in the worlds they hold for. A general repair, which
corrects a unit, applies a scaling law or a parameterization fitted over a wider range, or removes a defect,
contains no coefficient fitted to the test world and is built to hold on another world within the limits of its
laws; a general method holds wherever its table is computed from that world's own column, as for the curved
geometry of a tall atmosphere; and a calibration, fitted to a reference on columns or states of the test world, holds near them and is refitted
whenever the state or the world changes. Settings, added capabilities, couplings, a
bound, a mitigation and avoidances by configuration make up the remaining changes. The class of a repair followed
largely from the form of the Earth value it removed. Of the four forms, moreover, the fits had the largest effects on the test world's climate, effects of the order of the spreads measured among exoplanet GCMs, and every lunar climate number is therefore a conditional model result, holding given the
scenario and the Earth values that remain.

Between the two atmospheric models, the rain and the upper air still differ. Forced through the whole column with
the GCM's mean ascent, from its climate before and after the turn, CM1's boxes rain 75--84\% of the GCM's rain, little
of it at night over the sea, when the GCM rains half of its own there, while boxes that make their own ascent rain
about a tenth more than the GCM at a land site at 0\textdegree~E and twice its rain at sea.\claim{G21}\claim{G26}\claim{G28}
At those two sites the two ways of forcing the boxes bracket the GCM's rain, their spread wider than the difference
itself, whereas at a dry inland site both rain less; the boxes forced by the GCM's ascent also run 3.5--10~K colder
at 40--60~km, where the land box making its own ascent holds within 0.15~K of the GCM's air, and no box passes the
agreement test, the sea boxes reading as observations.\claim{G29} The ensemble's next
steps start there. With the boxes' ascent settled, the GCM's convection would be closed toward the free troposphere
such boxes support, PlaSim's cloud calibrated against the resolved cloud of each world, and the broadband fits
replaced by a radiation scheme reduced with existing tools from the reference model's spectroscopy
\autocite{hogan2022tool,hogan2020evaluating}, once that spectroscopy is checked against an independent code. Runs on planets that other studies model, among them slowly
rotating planets of \textcite{way2018climates} and the cases of the SAMOSA intercomparison
\autocite{haqqmisra2022samosa}, would then test how far the general repairs hold.
